%% A282_Die_Rechenkugel_En.tex
%% FFGFT A-Series -- A282: The Reckoning Bead
%% Autor: Johann Pascher  --  ORCID 0009-0000-6518-4064
%% Language: English  --  Engine: LuaLaTeX
%% Check script: a282_rechenkugel.py (checks 1--11)

\documentclass[11pt,a4paper]{article}

% ---- Engine: LuaLaTeX ----
\usepackage{fontspec}
\usepackage{mathtools}
\usepackage{unicode-math}
\usepackage{polyglossia}
\setmainlanguage{english}

\usepackage{microtype}
\emergencystretch=3.5em
\tolerance=2500
\hbadness=3000
\usepackage{enumitem}
\usepackage{booktabs}
\usepackage{longtable}
\usepackage{array}
\usepackage[a4paper,margin=2.6cm,headsep=1.1em,headheight=22pt]{geometry}
\usepackage{fancyhdr}
\usepackage{titlesec}
\usepackage{xcolor}
\usepackage{needspace}
\usepackage{tocloft}

% ---- Fonts ----
\setmainfont{Inter}[
  Scale=1.02,
  Ligatures=TeX ]
\setsansfont{Inter}[Scale=MatchLowercase,Ligatures=TeX]
\setmonofont{JetBrains Mono}[Scale=0.88]
\setmathfont{Libertinus Math}

% ---- Colours ----
\definecolor{accent}{HTML}{1F3A5F}
\definecolor{accent2}{HTML}{6A4E2E}
\definecolor{rulecol}{HTML}{B9C2CE}
\definecolor{linknav}{HTML}{2C5F8A}
\definecolor{linkcite}{HTML}{7A5C2E}

% ---- Schichtmarkierungen (A-Serie) ----
\definecolor{laySet}{HTML}{7A2E2E}
\definecolor{layK}{HTML}{1D5C3A}
\definecolor{layB}{HTML}{1F3A5F}
\definecolor{layS}{HTML}{6A4E2E}
\definecolor{layH}{HTML}{5A2E6A}

\newcommand{\LS}{\textcolor{laySet}{\textbf{\footnotesize\textsf{[STIPULATION]}}}\,}
\newcommand{\LK}{\textcolor{layK}{\textbf{\footnotesize\textsf{[K]}}}\,}
\newcommand{\LB}{\textcolor{layB}{\textbf{\footnotesize\textsf{[B]}}}\,}
\newcommand{\LSk}{\textcolor{layS}{\textbf{\footnotesize\textsf{[S]}}}\,}
\newcommand{\LH}{\textcolor{layH}{\textbf{\footnotesize\textsf{[H]}}}\,}

\newcommand{\pruef}[1]{{\small\itshape(check script, #1)}}

% ---- Hyperref ----
\usepackage[
  colorlinks=true,
  linkcolor=linknav,
  citecolor=linkcite,
  urlcolor=linknav,
  bookmarksnumbered=true,
  bookmarksopen=true,
  pdfstartview=FitH ]{hyperref}
\hypersetup{
  pdftitle={The Reckoning Bead (A282)},
  pdfauthor={Johann Pascher},
  pdfsubject={Token computation as the limiting case of Landauer accounting},
  pdfkeywords={Landauer, abacus, token, phase space, heat bath, FFGFT, A282}
}

% ---- Section styling ----
\titleformat{\section}
  {\color{accent}\normalfont\Large\bfseries}
  {\color{accent2}\thesection}{0.7em}{}
\titleformat{\subsection}
  {\color{accent}\normalfont\large\bfseries}
  {\color{accent2}\thesubsection}{0.6em}{}
\titlespacing*{\section}{0pt}{1.6\baselineskip}{0.7\baselineskip}
\titlespacing*{\subsection}{0pt}{1.1\baselineskip}{0.4\baselineskip}

% ---- Headers ----
\pagestyle{fancy}
\fancyhf{}
\renewcommand{\headrulewidth}{0.4pt}
\renewcommand{\headrule}{\hbox to\headwidth{\color{rulecol}\leaders\hrule height \headrulewidth\hfill}}
\fancyhead[L]{\footnotesize\color{accent}\itshape FFGFT $\cdot$ A282 $\cdot$ The Reckoning Bead}
\fancyhead[R]{\footnotesize\color{accent2}\thepage}

% ---- TOC look ----
\renewcommand{\cftsecfont}{\normalfont}
\renewcommand{\cftsecpagefont}{\normalfont}
\renewcommand{\cftsecleader}{\cftdotfill{\cftdotsep}}
\setlength{\cftbeforesecskip}{3pt}

\setlength{\parskip}{0.35em}
\setlength{\parindent}{0pt}
\linespread{1.06}

\begin{document}

% ===================== TITLE =====================
\begin{titlepage}
\centering
\vspace*{2.2cm}
{\color{rulecol}\rule{0.5\textwidth}{0.6pt}}\\[1.4cm]
{\color{accent}\Huge\bfseries The Reckoning Bead\par}
\vspace{1.0cm}
{\large\itshape Token computation as the limiting case\\ of Landauer accounting\par}
\vspace{0.6cm}
{\color{accent2}\rule{0.28\textwidth}{0.4pt}}\\[1.0cm]
{\large FFGFT A-Series $\cdot$ A282\par}
\vspace{0.4cm}
{\normalsize Third piece alongside A280 (accounting) and A281 (textual criticism)\par}
\vfill
{\large Johann Pascher\par}
\vspace{0.3cm}
{\footnotesize ORCID: 0009-0000-6518-4064\par}
\vspace{0.2cm}
{\footnotesize July 2026\par}
\vspace{1.4cm}
{\color{rulecol}\rule{0.5\textwidth}{0.6pt}}
\vspace*{1.5cm}
\end{titlepage}

% ===================== TOC =====================
\tableofcontents
\thispagestyle{fancy}
\clearpage

% ===================== ABSTRACT =====================
\begingroup\small\noindent
\textbf{Abstract.} Computing with tokens --- abacus beads, shells, coins --- is
the limiting case in which Landauer's argument comes apart into its two halves.
The \emph{accounting} holds there exactly and is for once visible: the stock from
which the tokens come and into which they flow back is a bath one can point at.
The \emph{thermal conversion} $\Delta S \to Q = T\Delta S$ does not hold there: an
abacus bead costs about $5\times10^{15}$ times the Landauer bound to move, its
positional bit carries $4\times10^{-23}$ of its own thermal entropy, and its
positional barrier stands at $3.6\times10^{15}\,kT$. The token computer therefore
falls squarely under the very restriction Landauer himself stated and the
reception dropped (A281): the counted states are not a thermal ensemble. From
this follows an ordering that runs against intuition: \textbf{visibility of the
structure and binding force of the number run in opposite directions.} The cruder
the carrier, the clearer the accounting and the more irrelevant the bound. The
passage back is continuous --- shrink the token until its barrier approaches
$kT$ and it becomes a Brownian reckoning bead, which is Bérut's colloid.
\par\endgroup

\bigskip

\begingroup\small\noindent
\textbf{Layer markings.} \LS declared starting point (German corpus token:
[SETZUNG]) --- \LB derived algebraically/mathematically from declared foundations
--- \LK empirically tested or documented from primary sources --- \LSk
plausibility sketch --- \LH open research question.
\par\endgroup

\bigskip

\begingroup\small\noindent
\textbf{Check script.} \texttt{a282\_rechenkugel.py} (checks 1--11). Every number
in this document is backed there by assertions; the parameter choices are marked
in the script as stipulations and can be varied freely.
\par\endgroup

\bigskip

% ===================== 1 =====================
\section{The stock as a visible bath}

A280 works with the world ocean: a glass is poured in, the volume does not vanish,
but nobody gets it back. The image is correct and has one drawback --- the bath
has to be inferred. Nobody sees the $10^{23}$ degrees of freedom into which the
bit volume tips.

\LS With token computation this is different. An abacus, a heap of shells, a till
with coins: next to the working area there lies a \textbf{stock}. The tokens come
from it and flow back into it. The stock is a bath one can point at.

\begin{center}
\begin{minipage}{0.80\textwidth}
\ttfamily\small
\begin{tabbing}
\ Working area\ \ \ \ \ \ \ \ \ \ \ \=\ Stock (bath)\\[0.3em]
+-----------------+\>+-------------------------+\\
|\ o\ o\ .\ o\ .\ .\ o\ .\ |\>|\ o o o o o o o o o o o o |\\
|\ .\ o\ o\ .\ o\ o\ .\ o\ |\ \ ---\textgreater\ \>|\ o o o o o o o o o o o o |\\
+-----------------+\>+-------------------------+\\
\ W = 2\^{}N config.\>\ tokens conserved, arrangement mixed
\end{tabbing}
\end{minipage}
\end{center}

\LB The structure is step for step the same as in A280:

\begin{itemize}[leftmargin=1.4em,itemsep=0.3em]
  \item \textbf{Region.} A token in the left position and the same token in the
        right position are two distinguishable macro-regions.
  \item \textbf{Region reduction.} Clearing the board maps all configurations onto
        one --- all tokens in the stock. Non-bijective, exactly like
        $\{A,B\}\to Z$.
  \item \textbf{Conservation.} No token vanishes. They spread into an arrangement
        from which the previous configuration can no longer be read off. That is
        the Liouville analogue --- and here it is literally countable.
  \item \textbf{Content neutrality.} The stock does not know whether a token stood
        for a debit, for a five, or for nothing. It absorbs it the same way.
\end{itemize}

\LB Cash with coins and a till is the same figure, and the double meaning of the
word “accounting” stops being a metaphor here: the balance of a till and
phase-space accounting have the same form. Barter is the case without tokens ---
there is no stock there and correspondingly no region reduction, only bijective
exchange.

% ===================== 2 =====================
\section{The accounting holds exactly}

\LB For $N$ tokens with two positions each, before clearing
$W_{\text{before}} = 2^N$ and after $W_{\text{after}} = 1$, hence
\begin{equation}
  \frac{\Delta S}{k_B} \;=\; \ln\frac{W_{\text{before}}}{W_{\text{after}}}
  \;=\; N \ln 2 .
  \label{eq:combinatorics}
\end{equation}

This is exact and carrier-independent \pruef{check 10}. It holds for beads, for
shells, for transistors, for colloids. Nothing in
equation~\eqref{eq:combinatorics} demands that the token be small, light, or
thermally driven. The count is combinatorics over regions --- exactly what A280
identifies as primary.

\LB Up to this point the token computer is a complete model of Landauer
accounting, and a better one than any other, because every step can be seen.

% ===================== 3 =====================
\section{The bound does not hold}

\LS For the following numbers it is stipulated: an abacus bead of mass
$m = 0.5$\,g, a displacement $d = 1$\,cm, a friction coefficient $\mu = 0.3$, an
ambient temperature $T = 300$\,K, a molar mass of $100$\,g/mol and a molar entropy
of $50$\,J/(mol$\cdot$K) for a solid. All stipulations are marked in the check
script and can be changed; the orders of magnitude do not depend on their precise
values.

\LB \textbf{The work per bead displacement.} Sliding friction gives
$W = \mu\,m\,g\,d = 1.47\times10^{-5}$\,J \pruef{check 2}. The Landauer bound at
300\,K is $k_B T \ln 2 = 2.87\times10^{-21}$\,J \pruef{check 1}. The ratio is
\begin{equation}
  \frac{W_{\text{bead}}}{k_B T \ln 2} \;\approx\; 5.1\times10^{15}
  \qquad \text{\pruef{check 3}} .
\end{equation}

\subsection{The scale band}

\begin{center}
\small
\begin{tabular}{@{}lll@{}}
\toprule
\textbf{Carrier} & \textbf{Cost per switching event} & \textbf{Multiple of $k_BT\ln 2$} \\
\midrule
colloid, quasistatic [B4][B5] & ${\sim}3\times10^{-21}$\,J & 1--10 \\
CMOS transistor (A280) & $10^{-17}$--$10^{-16}$\,J & $3.5\times10^{3}$ -- $3.5\times10^{4}$ \\
abacus bead & $1.5\times10^{-5}$\,J & $5.1\times10^{15}$ \\
\bottomrule
\end{tabular}
\end{center}

\LB Between bead and CMOS lie roughly \textbf{eleven decades} \pruef{check 5}. The
ordering is thus the reverse of the obvious conjecture: Landauer binds not at the
primitive end but at the small one. The cruder the carrier, the further the bound
lies below the actual expenditure.

\subsection{Why the bead pays nothing}

Two numbers account for the finding, and both concern the bead itself, not the
accounting.

\LB \textbf{First: the positional bit is vanishing against the bead's own
entropy.} A bead of 0.5\,g carries, under the stipulations above, an entropy of
$0.25$\,J/K, that is $1.8\times10^{22}\,k_B$ \pruef{check 6}. The positional bit
carries $\ln 2 = 0.69\,k_B$. Its share is
\begin{equation}
  \frac{\ln 2}{S_{\text{bead}}/k_B} \;\approx\; 3.8\times10^{-23}
  \qquad \text{\pruef{check 7}} .
\end{equation}
Moving the bead changes essentially nothing in the thermally accessible phase
space. The region reduction takes place in a coordinate decoupled from the bead's
$10^{22}$ thermal degrees of freedom.

\LB \textbf{Second: the positional barrier is thermally insurmountable.} It stands
at $3.6\times10^{15}\,k_BT$ \pruef{check 8}. An Arrhenius lifetime
$\propto \exp(3.6\times10^{15})$ is not numerically representable. There is no
thermal transfer between “left” and “right”, and therefore no equilibrium ensemble
over those two states.

\LB The bead is thus the extreme case of the row “emergent two-valuedness” in
A280: the two is a partition imposed on a quasi-continuum of $10^{22}$
microstates, and it survives not because it is fundamental but because the barrier
is absurdly high.

% ===================== 4 =====================
\section{The two halves of Landauer's argument}
\label{sec:twohalves}

\LB The token computer thereby does something no other model does as cleanly: it
separates the two steps of which Landauer's argument consists, and shows that they
are independent.

\begin{center}
\small
\begin{tabular}{@{}p{0.16\textwidth}p{0.36\textwidth}p{0.40\textwidth}@{}}
\toprule
\textbf{Step} & \textbf{Statement} & \textbf{Presupposition} \\
\midrule
combinatorics & $\Delta S = k_B \ln(W_{\text{b}}/W_{\text{a}})$ &
regions are distinguishable and countable \\
conversion & $Q = T\,\Delta S$ &
the counted states are thermally occupied \\
\bottomrule
\end{tabular}
\end{center}

\LB The first step holds exactly for the bead. The second does not hold for it at
all. The usual presentation runs the two together in one move and thereby makes
the second presupposition invisible. Only when a carrier violates it maximally
does it come into view.

\begin{quote}
\LB Landauer's bound is not the statement “region reduction costs”. It is the
statement “region reduction \emph{in thermally occupied states} costs”. The
qualifier is not cosmetic --- it carries the whole conversion.
\end{quote}

\LB This sharpens A280 at a place where it needs to be more precise. A280 §6
requires, for classical systems, distinguishability \emph{above} the thermal
resolution limit. That is the condition for two regions to count as two. It is
\emph{not} the condition for their reduction to cost heat --- for that one
additionally needs the barrier \emph{not} to lie arbitrarily far above $k_BT$.
The two conditions point in opposite directions, and between them lies the narrow
band in which Landauer binds.

% ===================== 5 =====================
\section{Landauer's own restriction}

\LK The finding is not a new discovery but the extreme case of a restriction
Landauer stated himself in 1961: his reset operation was applied to a thermal
equilibrium ensemble; for information not yet thermalised the calculation does not
apply directly [B1]. A281 documents the passage and shows that the reception
dropped it.

\LB With the abacus the restriction is maximally violated, and in two ways at
once:

\begin{enumerate}[leftmargin=1.6em,itemsep=0.3em]
  \item \textbf{No thermal ensemble.} The distribution over “left” and “right”
        does not arise from thermal occupation but from the hand that placed the
        bead.
  \item \textbf{Complete knowledge.} Whoever placed the beads knows the
        configuration. In the feedback bound (A281) this makes $I = \ln 2$ per
        token and
        \begin{equation}
          Q_{\min} \;=\; k_B T\,(\ln 2 - I) \;=\; 0
          \qquad \text{\pruef{check 11}} .
        \end{equation}
\end{enumerate}

\LB Thermodynamically, then, clearing a known abacus costs nothing. What it does
cost is friction --- and friction is an engineering quantity with no lower bound
in nature, reducible at will by better bearings.

% ===================== 6 =====================
\section{The passage back: the Brownian reckoning bead}

\LB The passage back into Landauer's domain of validity is continuous and can be
computed. Setting the positional barrier equal to $k_BT$ gives, for the critical
token mass,
\begin{equation}
  \mu\,m\,g\,d \;=\; k_B T
  \quad\Longrightarrow\quad
  m_{\text{crit}} \;=\; \frac{k_B T}{\mu\,g\,d} \;\approx\; 1.4\times10^{-19}\,\text{kg}
  \qquad \text{\pruef{check 9}} .
\end{equation}

\LK For comparison: a silica colloid of 2\,\textmu m diameter weighs
$8.4\times10^{-15}$\,kg, roughly $6\times10^{4}$ times more \pruef{check 9}. But
it holds its position not by friction, rather by an optical trap, and the height
of that barrier is adjustable down to a few $k_BT$. That is exactly what makes it
a Brownian reckoning bead --- and exactly why Bérut's apparatus [B4] is the only
one in which the bound actually binds.

\begin{quote}
\LB Bérut's colloid is not a model of the bit. It is an abacus bead small enough
for the bath to push it.
\end{quote}

\LB The series thereby closes: shell, coin, abacus bead, transistor, colloid. It
is a scale, not a jump. What changes along the scale is not the accounting --- it
is the same everywhere --- but the ratio of the positional barrier to $k_BT$.

% ===================== 7 =====================
\section{Two axes, running opposite}

\LB From this follows the ordering that gives the document its yield:

\begin{center}
\small
\begin{tabular}{@{}lll@{}}
\toprule
\textbf{Carrier} & \textbf{Visibility of the structure} & \textbf{Binding force of the bound} \\
\midrule
abacus, shells, coins & complete (bath can be pointed at) & none ($10^{15}$ above) \\
CMOS & slight (bath is the lattice) & weak ($10^{4}$ above) \\
colloid in an optical trap & none (bath not separable) & full (factor 1--10) \\
\bottomrule
\end{tabular}
\end{center}

\begin{quote}
\LB At the most primitive level the accounting is maximally visible and the bound
maximally irrelevant. Visibility of the structure and binding force of the number
are two different axes, and they run in opposite directions.
\end{quote}

\LB This also explains why intuition regularly misleads in this matter. Anyone
wanting to make Landauer plausible reaches for the coarse image --- the bead, the
drawer, the glass of water --- and thereby lands exactly where the bound has
ceased to mean anything. The image carries the structure and loses the number.

% ===================== 8 =====================
\section{The token stock as a cost category of its own}

\LB One point remains that thermodynamic accounting does not capture and that
becomes especially clear with the token computer: \textbf{without tokens one
cannot compute.} The stock is an exhaustible resource. An abacus with too few
beads does not compute more slowly; it does not compute at all.

\LB This is a \emph{material bound}, not a thermodynamic one. It has no lower
limit in $k_BT$; it has a lower limit in the number of items. It is related to
what A280 §7.3 opens up --- the computer pays for existing, not only for computing
--- but it is not the same thing: there the subject is a running entropy current,
here a standing inventory.

\LSk With semiconductors this category drops out of view, because the stock is
cast in silicon: the regions are laid down in fabrication and not drawn during
operation. The inventory is there, so it is not counted. On the abacus it is
visible, because it can be used up.

\LH \textbf{Open question.} Whether inventory and current can be carried in a
single ledger --- token stock and entropy current as two balances of one computer
--- is open. The question is the same as the one at the end of A281, posed from
the material side.

% ===================== 9 =====================
\section{Placement in the corpus}

\LB The three pieces stand to one another as follows:

\begin{itemize}[leftmargin=1.4em,itemsep=0.3em]
  \item \textbf{A280} does the accounting: region, region reduction, Liouville,
        bath, bound. The physical core.
  \item \textbf{A281} does the textual criticism: what Landauer proved in 1961,
        his own ensemble restriction, the reception history, the language.
  \item \textbf{A282} exhibits the limiting case at which both findings converge
        and become visible: the token computer satisfies A280's accounting exactly
        and violates A281's ensemble condition maximally.
\end{itemize}

\LB The dictionary of A281 §1.1 applies unchanged. In the language of this
document: token $=$ carrier $=$ region; stock $=$ bath; clearing the board $=$
carrier operation $=$ region reduction.

\LB \textbf{Contact with FFGFT: none.} Like A280 and A281, this document lives at
the level of statistical mechanics; the $\xi$ scheme is untouched.

% ===================== 10 =====================
\section{Layer-status summary}

\begin{center}
\small
\begin{longtable}{@{}p{0.72\textwidth}l@{}}
\toprule
\textbf{Statement} & \textbf{Status} \\
\midrule
\endfirsthead
\toprule
\textbf{Statement} & \textbf{Status} \\
\midrule
\endhead
The token stock is a bath analogue one can point at & \LS \\
Parameters of the abacus bead (mass, distance, friction, $T$) & \LS \\
Clearing the board is a non-bijective region map & \LB \\
Tokens are conserved (Liouville analogue) & \LB \\
$\Delta S/k_B = N\ln 2$, exact and carrier-independent & \LB \\
$k_BT\ln 2 = 2.87\times10^{-21}$\,J at 300\,K & \LK \\
Bead displacement $1.47\times10^{-5}$\,J & \LB \\
Bead $\approx 5.1\times10^{15}\times$ Landauer & \LB \\
Scale band colloid / CMOS / bead & \LK \\
Eleven decades between bead and CMOS & \LB \\
Bead's own entropy $1.8\times10^{22}\,k_B$ & \LB \\
Positional bit $= 3.8\times10^{-23}$ of that entropy & \LB \\
Positional barrier $3.6\times10^{15}\,k_BT$ & \LB \\
Combinatorics and thermal conversion are independent & \LB \\
The conversion presupposes thermally occupied states & \LB \\
Landauer's ensemble restriction (original passage) & \LK \\
Known abacus: $I=\ln 2$, hence $Q_{\min}=0$ & \LB \\
Critical token mass $1.4\times10^{-19}$\,kg & \LB \\
Silica colloid 2\,\textmu m: $8.4\times10^{-15}$\,kg & \LK \\
Bérut's colloid as a Brownian reckoning bead & \LB \\
Visibility and binding force run in opposite directions & \LB \\
Token stock as a material bound & \LB \\
Why the category drops out of view with semiconductors & \LSk \\
A single ledger for inventory and current & \LH \\
\bottomrule
\end{longtable}
\end{center}

\textbf{Balance:} 15$\times$ \LB $\cdot$ 4$\times$ \LK $\cdot$ 2$\times$ \LS
$\cdot$ 1$\times$ \LSk $\cdot$ 1$\times$ \LH $=$ 23 entries.

% ===================== SOURCES =====================
\clearpage
\phantomsection
\addcontentsline{toc}{section}{Sources}
\section*{Sources}

\begingroup
\small
\begin{description}[leftmargin=3.2em,style=sameline,itemsep=0.3em]

\item[{[B1]}] R. Landauer, \emph{Irreversibility and Heat Generation in the
Computing Process}, IBM J.\ Res.\ Dev.\ \textbf{5}, 183 (1961).

\item[{[B2]}] Liouville's theorem: standard result of Hamiltonian mechanics;
e.g.\ Landau/Lifshitz, \emph{Mechanics}, §46.

\item[{[B3]}] C.~H. Bennett, \emph{The Thermodynamics of Computation --- a
Review}, Int.\ J.\ Theor.\ Phys.\ \textbf{21}, 905 (1982).

\item[{[B4]}] A. Bérut et al., \emph{Experimental verification of Landauer's
principle}, Nature \textbf{483}, 187 (2012).

\item[{[B5]}] Y. Jun, M. Gavrilov, J. Bechhoefer, \emph{High-Precision Test of
Landauer's Principle in a Feedback Trap}, Phys.\ Rev.\ Lett.\ \textbf{113},
190601 (2014).

\item[{[B6]}] T. Sagawa and M. Ueda, \emph{Nonequilibrium Thermodynamics of
Feedback Control}, Phys.\ Rev.\ E \textbf{85}, 021104 (2012).

\item[{[B7]}] Semiclassical state counting $\Gamma/h^f$ and the entropy of
solids: standard; e.g.\ Landau/Lifshitz, \emph{Statistical Physics}, §7 and §64.

\item[{[B8]}] CODATA 2018 values; $k_B = 1.380649\times10^{-23}$\,J/K exactly
defined since the SI revision of 2019.

\end{description}
\endgroup

\vspace{1em}
\noindent
{\small\itshape The numerical values in this document come from the check script
\texttt{a282\_rechenkugel.py}; the parameters set there are model choices, not
measured values. The orders of magnitude are insensitive to their variation.}

\end{document}
